\section{Papal Service: France, the Sixtine Index, the Vulgate, the
Catechisms (1590--1599)}

\subsection{The French legation of 1590}

In the autumn of 1589 the assassination of Henry~III had thrown
France into the last convulsion of its religious wars. ``In 1590,''
the 1907 Catholic Encyclopedia records, ``he went with Cardinal
Gaetano as theologian to the embassy Sixtus~V was then sending into
France to protect the interests of the Church.'' The autobiography
turns the legation into lived history: Bellarmine entered Paris with
the legation in January 1590 (so the editors' notes date it), was
shut in through Henry of Navarre's siege of the city that summer,
and remembered---with the taste for private prophecy that runs
through the memoir---that when letters from Rome
reached the legate in early September, everyone feared their
contents, for ``Pope Sixtus was angered at the Cardinal, at his
secretary, and at Bellarmine himself, because of the proposition
found in his books denying that the Pope is the direct lord of the
whole world''; and that ``N.'' said, ``In these letters is contained
the death of Pope Sixtus~V''---and so, to universal astonishment, it
proved (D\"ollinger--Reusch, 37--38; A, retrospective
self-presentation; the prophecy belongs to the tradition audit, not
to history). Sixtus had died on 27~August 1590. Bellarmine, gravely
ill with dysentery on the road home, was carried by litter and
reached Rome on 11~November 1590, passing unrecognized, as he
noted, through Basel (A).

\subsection{The Sixtine Index episode, precisely stated}

The anger behind those letters is the most instructive censure
story of the age, and it must be told exactly. In \work{De Summo
Pontifice} (book~V) Bellarmine had rejected the opinion that the
pope possesses \emph{direct} temporal dominion over the whole
world, teaching instead an \emph{indirect} power in temporal
matters, ordered to spiritual ends. Sixtus~V held the stronger
view. The autobiography states what the pope did: ``on account of
that proposition about the Pope's direct dominion over the whole
globe, Sixtus placed his \work{Controversies} on the Index of
prohibited books, \latin{donec corrigerentur} [until they should be
corrected]; but when he died, the Sacred Congregation ordered his
name to be deleted from the book of the Index''
(D\"ollinger--Reusch, 38--39; A). D\"ollinger and Reusch, from the
bibliography of the Index, add the decisive bibliographical facts:
the listing stood in ``the very rare Index of 1590,'' with the
formula \latin{donec corrigantur}, and was struck from the
subsequent Indexes (K). The 1590 Sixtine Index was printed but
never effectively promulgated; Sixtus died before publication, and
his successors suppressed the whole edition. The 1907 Catholic
Encyclopedia completes the reversal: the threat ``was averted by
Sixtus's death,'' and Gregory~XIV gave the \work{Disputationes} a
special approbation (E/K, reported).

\witnesslimit{Three claims must not be run together: (1) that
Sixtus~V placed the first volume on his 1590 Index pending
correction is secure (A, with K bibliographical confirmation); (2)
that the 1590 Index never acquired binding force is the standard
bibliographical judgment, reported here from
D\"ollinger--Reusch's Index scholarship; (3) any statement about
Sixtus's motives beyond the temporal-power proposition, or about
what ``correction'' would have been demanded, is interpretation.
Bellarmine's own narrative is retrospective and far from neutral:
he tells the story as the prelude to his repayment of ``good for
evil'' in the Vulgate affair.}

The episode left a permanent irony that the rest of the life kept
sharpening: the theologian too papal for Paris was, in 1590, too
little papal for the pope. Twenty years later the Parlement of
Paris would condemn his defense of the indirect power as an attack
on the crown (\S\ref{sec:penwars}); the via media satisfied
neither extreme, a fact the 1907 Catholic Encyclopedia already
framed as the signature of his position.

\subsection{The Vulgate: repaying good for evil}

Trent had called for an authoritative edition of the Latin Bible;
Sixtus~V had produced one in 1590, largely by his own editorial
hand, and it was full of unfortunate readings---in the
autobiography's blunt phrase, ``very many things wrongly changed''
(\latin{permulta perperam mutata}). On Sixtus's death the question
was what to do with it. Grave men counseled public prohibition---a
staggering step against a pope's own Bible. Bellarmine, before
Gregory~XIV, argued the other course: the Sixtine Bible should not
be prohibited but corrected, reissued under Sixtus's own name, with
a preface explaining that errors---of the printers or
otherwise---had crept into the first edition through haste
(D\"ollinger--Reusch, 38--39; A). That is what was done. A small
commission worked at Zagarolo in the house of Marco Antonio
Colonna---Bellarmine names Cardinal Colonna, the English Cardinal
Allen, the Master of the Sacred Palace, and himself among some
seven or eight members---and the corrected Bible appeared under
Clement~VIII in 1592, still bearing Sixtus's name, with a preface
``which the same N.\ composed'' (A). The 1907 Catholic
Encyclopedia confirms: ``Bellarmine himself wrote the preface,
still prefixed to the Clementine edition.''

Bellarmine presents the episode as returning ``good for evil'' to
the pope who had put him on the Index. Critics from the
seventeenth century onward have read the preface's account of
printers' errors as a diplomatic fiction protecting Sixtus's
memory---the charge of \latin{pia fraus} was urged against him in
his own beatification process, as the eighteenth-century vota
show. This study records both the document and the debate: the
preface's strategy is Bellarmine's own testimony; its moral
character was and remains argued (unresolved; see the tradition
audit on the cause).

\subsection{The catechisms}

Recalled by Clement~VIII to Rome after about two years in Naples
(c.~1596--97) as papal theologian and consultor of the Holy Office
(\work{Lux illa}, AAS~22, 595; M),
Bellarmine received from the pope a commission of an entirely
different scale of audience: a catechism for the unlettered. The
\work{Dottrina cristiana breve} appeared in 1597, and the
\work{Dichiarazione pi\`u copiosa della dottrina cristiana}---the
teacher's expansion---in 1598 (CE 1907; Benedict~XVI 2011). The
doctoral letter of 1931 gives the official measure of their
career: composed ``by order of Clement~VIII,'' the catechism set
out Catholic truth ``plainly, exactly, and in order,'' and ``for
almost three centuries'' fed the Christian people ``in many
regions of Europe and of the world''
(\work{Providentissimus Deus}, AAS~23, 436; M). Even as cardinal
and archbishop he kept teaching it to children and the unlettered
with his own mouth---a habit \work{Lux illa} singles out
(AAS~22, 595; M). The pairing is the man in miniature: the same
years that produced the folio artillery of the \work{Disputationes}
produced the two little books by which, for three centuries,
Italian children learned the sign of the cross.
